% Original comparison fixture, MIT, satz contributors (2026).
% KOMA-Script reference for the explicitly documented satz header policy.
\documentclass[11pt,twoside,open=any]{scrreprt}
\usepackage[T1]{fontenc}
\usepackage{lmodern}
\usepackage[a4paper,inner=25mm,outer=25mm,top=25mm,bottom=25mm]{geometry}
\usepackage{booktabs}
\usepackage{scrlayer-scrpage}
\usepackage[hidelinks,hypertexnames=false]{hyperref}
\automark[section]{chapter}
\clearpairofpagestyles
\lehead{\leftmark}
\rohead{\rightmark}
\lefoot*{\pagemark}
\rofoot*{\pagemark}
\renewcommand*{\chaptermarkformat}{}
\renewcommand*{\sectionmarkformat}{}
% Chapter marks provide the fallback until the chapter has its own section.
\renewcommand*{\chaptermark}[1]{\markboth{#1}{#1}}
\renewcommand*{\sectionmark}[1]{\markright{#1}}
% Preserve satz's existing raw-counter parity and displayed offset.
\renewcommand*{\thepage}{\number\numexpr\value{page}-1\relax}
\newcommand{\CaseText}[1]{%
  \textbf{#1}\par\medskip
  This report records a fictional instrument check. The values illustrate
  document structure and do not describe an experiment. Each page has a fixed
  break so that the LaTeX and Typst versions exercise the same header case.
  \par\medskip
  A useful running head identifies the current chapter or section without
  carrying information from a previous chapter. Chapter-opening pages leave
  the running head empty. Continuation pages retain their context even when
  the chapter has no sections of its own.\par
}
\begin{document}
\begingroup
\pagestyle{empty}
\renewcommand*{\chapterpagestyle}{empty}
\tableofcontents
\clearpage
\endgroup
\setcounter{page}{2}
\chapter{Instrument checks}\label{ch:instrument}
\section{Aim}
\CaseText{CASE opening-a}
\newpage

\section{Measurements}
\CaseText{CASE section-change}
\begin{table}[t]
\centering
\begin{tabular}{lr}
\toprule
\textbf{Setting} & \textbf{Example value}\\
\midrule
Reference & 1.00\\
Repeat & 1.02\\
\bottomrule
\end{tabular}
\caption{Illustrative readings; no experimental data.}\label{tab:readings}
\end{table}
\newpage

\CaseText{CASE continuation-a}
The readings in Table~\ref{tab:readings} belong to Chapter~\ref{ch:instrument}.
\begin{figure}[t]
\centering
\fbox{\rule{0pt}{32mm}\rule{.75\linewidth}{0pt}}
\caption{Placeholder for an instrument diagram.}\label{fig:instrument}
\end{figure}
\clearpage

\chapter{Section-free chapter}
\CaseText{CASE opening-b}
\newpage
\CaseText{CASE continuation-b-even}
\newpage
\CaseText{CASE continuation-b-odd}
\newpage
\section{Follow-up}
\CaseText{CASE follow-up-even}
\newpage
\CaseText{CASE follow-up-odd}
The diagram in Figure~\ref{fig:instrument} remains linked across chapters.
\clearpage

\addchap{Closing remarks}
\markboth{Closing remarks}{Closing remarks}
\CaseText{CASE opening-unnumbered}
\newpage
\CaseText{CASE unnumbered-odd}
\newpage
\setcounter{page}{4}
\CaseText{CASE counter-reset}
\end{document}
